\documentclass[11pt,a4paper]{article}
\usepackage[margin=2.3cm]{geometry}
\usepackage{amsmath,amssymb,amsfonts,bm}
\usepackage{booktabs,array,longtable}
\usepackage{graphicx}
\usepackage{xcolor}
\usepackage[colorlinks=true,linkcolor=blue!50!black,citecolor=blue!50!black,urlcolor=blue!50!black]{hyperref}
\usepackage{algorithm}
\usepackage{algpseudocode}

\newcommand{\vect}[1]{\bm{#1}}
\newcommand{\mat}[1]{\mathbf{#1}}
\newcommand{\skw}[1]{[\,#1\,]_{\times}}
\newcommand{\dd}{\,\mathrm{d}}
\newcommand{\softplus}{\operatorname{softplus}}
\newcommand{\relu}{\operatorname{relu}}

\title{Differentiable explicit impact model of visco-hyperelastic beam lattices with frictional contact\\[4pt]
\large Physics, contact, material, element, time integration, test rig and exact sensitivities}
\author{Model documentation\\[2pt]\small code: \texttt{src/explicitFE.py}, \texttt{explicitSolve.py}, \texttt{contact.py}, \texttt{latticeSetup.py},\\
\small \texttt{material.py}, \texttt{adjoint.py}, \texttt{objectives.py}, \texttt{optimizeDrop.py}}
\date{October 2026}

\begin{document}
\maketitle

\begin{abstract}
This document describes the current simulation model used to compute the impact response of additively
manufactured beam lattices and its exact design sensitivities. A lattice of slender struts made of an
elastomeric polyurethane (Carbon EPU~45) is bonded to a rigid base plate and capped by a rigid top plate; a
rigid impactor carrying a prescribed kinetic energy loads the top plate. The struts are modelled with a
corotational three-dimensional Timoshenko beam element whose section law is visco-hyperelastic (Yeoh
hyperelasticity integrated over the cross-section, quasi-linear Prony viscoelasticity). Contact between struts
is represented by a line-contact density integrated along the beams, contact with both plates by Hertzian node
contact, and the impactor--plate interaction by a critically damped Hertz law; regularised Coulomb friction acts
at every contact. The equations of motion are integrated with an explicit central-difference (velocity Verlet)
scheme with an exact exponential-map rotation update and an adaptive stable time step. Design gradients are
obtained with a checkpointed discrete adjoint of the time-stepping scheme. All model components are written in
PyTorch, so every force is differentiable with respect to positions, rotations and section properties.
\end{abstract}

\tableofcontents
\newpage

%=====================================================================================================
\section{Units, notation and overview}
\label{sec:units}
All quantities use the consistent set \textbf{mm, N, s, tonne (t), MPa}: $1\,\mathrm{t\,mm\,s^{-2}} = 1$\,N,
$1$\,J $=1000$\,N\,mm and a density of $1\,\mathrm{kg\,m^{-3}} = 10^{-12}\,\mathrm{t\,mm^{-3}}$.
Vectors are bold, $\skw{\vect{a}}$ is the skew matrix with $\skw{\vect{a}}\vect{b} = \vect{a}\times\vect{b}$,
and $\exp(\skw{\vect{\theta}})$ is the rotation by the rotation vector $\vect{\theta}$.
$\softplus(s)=\ln(1+e^{s})$, $\sigma(s) = 1/(1+e^{-s})$ and $\relu(s)=\max(s,0)$.

The model has the following parts, described in turn:
\begin{enumerate}
\item material (Section~\ref{sec:material}): EPU~45 visco-hyperelastic constants;
\item lattice geometry, sections and design variables (Section~\ref{sec:geometry});
\item corotational Timoshenko beam element (Section~\ref{sec:element});
\item section law: fibre-integrated Yeoh hyperelasticity with Prony viscoelasticity (Section~\ref{sec:section});
\item inertia (Section~\ref{sec:mass});
\item test rig and boundary conditions: bonded base plate, tied top plate, impactor (Section~\ref{sec:rig});
\item contact: impactor--plate, beam--beam, beam--plate, friction (Section~\ref{sec:contact});
\item explicit time integration and the stable time step (Section~\ref{sec:time});
\item energy accounting (Section~\ref{sec:energy});
\item response measures, objectives and exact sensitivities (Section~\ref{sec:sens});
\item verification (Section~\ref{sec:verification}) and limitations (Section~\ref{sec:limits}).
\end{enumerate}

%=====================================================================================================
\section{Material: Carbon EPU 45}
\label{sec:material}
EPU~45 is an elastomeric polyurethane printed by digital light synthesis (DLS). Its response is strongly
rate-dependent and dissipative. The model uses a hyperelastic long-term response with a single relaxation term.

\subsection{Constants}
\begin{table}[h]
\centering
\small
\begin{tabular}{p{0.36\textwidth}p{0.33\textwidth}p{0.24\textwidth}}
\toprule
quantity & value & note\\
\midrule
density $\rho$ & $1.06\times10^{-9}\,\mathrm{t\,mm^{-3}}$ ($1060\,\mathrm{kg\,m^{-3}}$) & Carbon TDS (ASTM D792)\\
Yeoh $C_{10}, C_{20}, C_{30}$ & $2.4,\ 0.15,\ 0.02$\,MPa & long-term, compression-calibrated\\
compressibility $D_i$ & $\approx 0$ & nearly incompressible\\
Prony $g_1$, $k_1$ & $0.958$, $0$ & shear relaxation only\\
Prony $\tau_1$ & $5.2\times 10^{-5}$\,s & \\
Poisson ratio $\nu$ & $0.49$ & assumed\\
long-term modulus $E_\infty = 6C_{10}$ & $14.4$\,MPa & small-strain limit of Yeoh\\
instantaneous modulus $E_0 = E_\infty/(1-g_1)$ & $342.9$\,MPa & governs the explicit time step\\
shear moduli $G_\infty$, $G_0$ & $E/(2(1+\nu))$ & \\
\bottomrule
\end{tabular}
\caption{EPU 45 constants (\texttt{config/materials.yaml}, entry \texttt{EPU\_45}; \texttt{material.epu45()}).}
\label{tab:material}
\end{table}

$E_\infty$ agrees with the manufacturer's tensile modulus (17\,MPa at 500\,mm/min, i.e.\ close to the relaxed
state). The Yeoh constants are calibrated in compression; they are stiffer than the tensile data sheet at large
tensile strain (5.5/11.3/40.8\,MPa vs.\ 3/4/9\,MPa at 50/100/200\,\% strain), which is the expected
tension/compression asymmetry of such fits. With one Prony term the storage and loss moduli are
\begin{equation}
E'(\omega) = E_\infty + (E_0-E_\infty)\frac{\omega^2\tau_1^2}{1+\omega^2\tau_1^2},\qquad
E''(\omega) = (E_0-E_\infty)\frac{\omega\tau_1}{1+\omega^2\tau_1^2},
\end{equation}
giving $\tan\delta = E''/E' = 0.73$ at 100\,Hz and $2.1$ at 1\,kHz, i.e.\ strong damping at impact
frequencies. (At 1\,Hz the model gives $E'=14.4$\,MPa, whereas the manufacturer's DMA reports
$\approx 65$\,MPa near $T_g=30\,^\circ$C; the slow viscoelasticity close to $T_g$ is not represented.)

\subsection{Yeoh hyperelasticity}
For an incompressible Yeoh solid the strain-energy density is
\begin{equation}
W(\bar I_1) = C_{10}(\bar I_1-3) + C_{20}(\bar I_1-3)^2 + C_{30}(\bar I_1-3)^3 .
\end{equation}
For uniaxial stretch $\lambda$, $\bar I_1 - 3 = \lambda^2 + 2/\lambda - 3$ and the nominal (first Piola) stress is
\begin{equation}
P(\lambda) = 2\left(\lambda - \lambda^{-2}\right)\left[C_{10} + 2C_{20}(\bar I_1-3) + 3C_{30}(\bar I_1-3)^2\right],
\qquad \left.\frac{\dd P}{\dd\lambda}\right|_{\lambda=1} = 6C_{10} = E_\infty .
\label{eq:yeoh-uniaxial}
\end{equation}
$P\to-\infty$ as $\lambda\to 0$, so a strut element cannot be compressed to zero length.
For simple shear $\gamma$, $\bar I_1 - 3 = \gamma^2$ and $\tau = 2\gamma\left[C_{10}+2C_{20}\gamma^2+3C_{30}\gamma^4\right]$.

%=====================================================================================================
\section{Lattice, cross-sections and design variables}
\label{sec:geometry}
\subsection{Lattice}
The lattice is the C4 cubic cell of the original optimisation study (exampleCase~3), $3\times3\times1$ cells of
size $12.7$\,mm, i.e.\ a $38.1\times38.1\times12.7$\,mm coupon with 168 nodes and 276 struts (strut lengths
2.6--6.4\,mm). Each strut is subdivided into elements of target length $h_e$ (number of elements per strut
$=\operatorname{round}(L/h_e)$); $h_e=1$\,mm (1332 elements) is the converged results mesh, $h_e=3$\,mm
(444 elements) is used for fast tests. The mapping element$\to$strut is kept (\texttt{beam\_of\_elem}).

\subsection{Circular sections}
Each element has a solid circular section of radius $r$:
\begin{equation}
A = \pi r^2,\quad A_{s} = \kappa A,\quad I_y = I_z = \tfrac{\pi}{4}r^4,\quad J = I_y + I_z,\qquad
\kappa = \frac{6(1+\nu)}{7+6\nu} = 0.899 \ \text{(Cowper)} .
\end{equation}
Struts are short and stout ($L/r = 3$--$8$), so the shear coefficient matters.

\subsection{Design variables and symmetry}
Design rows $\vect{x}\in[0,1]^{n}$ are mapped to radii by $r = r_{\min} + (r_{\max}-r_{\min})\,x$ inside the
objective (e.g.\ $[r_{\min}, r_{\max}] = [0.4, 1.2]$\,mm), then expanded by a symmetry rule
(\emph{Independent}: one row per strut; \emph{Unitcell}: one row per strut of the unit cell, repeated;
\emph{OneValue}: a single radius; \emph{Midplane}: three groups by height), and finally to every element of
each strut. The expansion is linear, so its derivative is exact.

%=====================================================================================================
\section{Corotational Timoshenko beam element}
\label{sec:element}
\subsection{Kinematic state}
Each node carries its position $\vect{x}_a = \vect{X}_a + \vect{u}_a$ and a rotation $\mat{R}_a\in SO(3)$
measured from the reference configuration. No element-level state is stored.

For element $e=(1,2)$ with reference axes $\mat{T}^0 = [\vect{e}^0_1, \vect{e}^0_2, \vect{e}^0_3]^\top$
(axis 1 along the chord; axis 2 from the global axis least aligned with the chord), the current nodal triads are
$\mat{t}_a = \mat{R}_a \mat{T}^{0\top}$ (columns $\vect{n}^a_x, \vect{n}^a_y, \vect{n}^a_z$).

\subsection{Element frame}
\begin{equation}
L = \lVert \vect{x}_2-\vect{x}_1\rVert,\quad
\vect{e}_1 = \frac{\vect{x}_2-\vect{x}_1}{L},\quad
\vect{q} = \tfrac12\left(\vect{n}^1_y + \vect{n}^2_y\right),\quad
\vect{e}_3 = \frac{\vect{e}_1\times\vect{q}}{\lVert \vect{e}_1\times\vect{q}\rVert},\quad
\vect{e}_2 = \vect{e}_3\times\vect{e}_1 .
\end{equation}
This frame follows Battini and Pacoste~\cite{battini2002} and needs no stored state.

\subsection{Local deformations}
The rotation of node $a$ relative to the element frame, $\mat{Q}_a = \mat{E}^\top\mat{t}_a$ with
$\mat{E} = [\vect{e}_1,\vect{e}_2,\vect{e}_3]$, is converted to a rotation vector by the logarithm
\begin{equation}
\vect{a} = \tfrac12\operatorname{vec}\left(\mat{Q}_a-\mat{Q}_a^\top\right),\quad s = \lVert\vect{a}\rVert = \sin\phi,
\quad c = \tfrac12(\operatorname{tr}\mat{Q}_a - 1) = \cos\phi,\quad
\vect{\theta}_a = \vect{a}\,\frac{\operatorname{atan2}(s,c)}{s}.
\label{eq:log}
\end{equation}
This form is smooth for $\phi\in[0,\pi)$. (The arcsine form used in the earlier linear code stiffens like
$1/\cos\phi$ and diverges at $90^\circ$ element-relative rotation, which occurs during crushing.)
The six local deformations are
\begin{equation}
\vect{d} = \left(\Delta L,\ r_x,\ r_{y1},\ r_{z1},\ r_{y2},\ r_{z2}\right),\quad
\Delta L = L - L_0,\quad r_x = \theta^{2}_x-\theta^{1}_x,\quad r_{ya}=\theta^{a}_y,\ r_{za}=\theta^{a}_z .
\end{equation}

\subsection{Internal force by virtual work}
With local resultants $\vect{f}_\mathrm{loc}(\vect{d}) = (N, T, M_{y1}, M_{z1}, M_{y2}, M_{z2})$ from the section
law (Section~\ref{sec:section}), the global nodal forces and moments are the exact virtual-work transpose
\begin{equation}
\begin{bmatrix}\vect{f}_u\\ \vect{f}_\theta\end{bmatrix}
= \left(\frac{\partial\vect{d}}{\partial(\vect{u},\delta\vect{\theta})}\right)^{\!\top}\vect{f}_\mathrm{loc},
\qquad \mat{R}_a \leftarrow \left(\mat{I} + \skw{\delta\vect{\theta}_a}\right)\mat{R}_a ,
\label{eq:fint}
\end{equation}
i.e.\ a vector--Jacobian product computed by automatic differentiation, with rotational variations taken as
spatial spins. For an elastic law, $\vect{f}_\mathrm{loc} = \partial\Pi/\partial\vect{d}$ and
(\ref{eq:fint}) is the exact gradient of the strain energy, so the explicit dynamics conserves energy. The
geometric (P--$\delta$) stiffness is included automatically.

\subsection{Linear Timoshenko resultants}
The small-strain part of the section law is the exact two-node Timoshenko element (per bending plane, with
$\Omega = EI/(G A_s L_0^2)$, $\mu = 1+12\Omega$, $\lambda=1+3\Omega$, $\gamma = 1-6\Omega$):
\begin{equation}
N = \frac{EA}{L_0}\Delta L,\quad T = \frac{GJ}{L_0}r_x,\quad
\begin{bmatrix}M_1\\M_2\end{bmatrix} = \frac{2EI}{\mu L_0}
\begin{bmatrix}2\lambda & \gamma\\ \gamma & 2\lambda\end{bmatrix}
\begin{bmatrix}r_1\\r_2\end{bmatrix} .
\label{eq:timo}
\end{equation}
Its small-displacement tangent equals the classical 12$\times$12 Timoshenko stiffness to $10^{-9}$.

%=====================================================================================================
\section{Section law: fibre-integrated Yeoh with Prony viscoelasticity}
\label{sec:section}
\subsection{Large-strain long-term response}
Axial stretch, bending and torsion all use the Yeoh law, integrated over the cross-section. From the local
deformations the element strains are
\begin{equation}
\varepsilon_0 = \frac{\Delta L}{L_0},\quad \kappa_y = \frac{r_{y2}-r_{y1}}{L_0},\quad
\kappa_z = \frac{r_{z2}-r_{z1}}{L_0},\quad \vartheta = \frac{r_x}{L_0},
\end{equation}
and a fibre at section coordinates $(y,z)$ has stretch and torsional shear
\begin{equation}
\lambda_f = 1 + \varepsilon_0 - y\kappa_z + z\kappa_y,\qquad \gamma_f = \vartheta\sqrt{y^2+z^2},\qquad
\bar I_1 - 3 = \lambda_f^2 + 2/\lambda_f - 3 + \gamma_f^2 .
\end{equation}
The element energy is
\begin{equation}
U_\mathrm{fib}(\vect{d}) = L_0 \int_{A} W\!\left(\bar I_1\right)\dd A \approx
L_0\, r^2 \sum_{k} w_k W\!\left(\bar I_1(y_k,z_k)\right),
\end{equation}
evaluated with a polar quadrature: 4 Gauss--Legendre rings ($\int_0^1 f(s)\,s\dd s$) $\times$ 16 angles.
$U_\mathrm{quad}$ is the same integral with the small-strain limit $W\to C_{10}(3\varepsilon^2+\gamma^2)$.
The long-term resultants are
\begin{equation}
\vect{f}_\infty(\vect{d}) = \vect{f}_\mathrm{lin}(\vect{d}) +
\frac{\partial}{\partial\vect{d}}\Big(U_\mathrm{fib}(\vect{d}) - U_\mathrm{quad}(\vect{d})\Big),
\label{eq:finf}
\end{equation}
with $\vect{f}_\mathrm{lin}$ from (\ref{eq:timo}) using $E_\infty, G_\infty$. At small strain the correction
vanishes and the validated linear Timoshenko element is recovered (including its shear flexibility); at large
strain axial, bending and torsional responses follow Yeoh, including the net axial force that pure bending
produces because of the tension/compression asymmetry. Transverse shear remains linear. The law is
conservative. This is the corotational fibre approach (e.g.\ \cite{battini2002}) with a hyperelastic fibre law;
see also geometrically exact beams~\cite{simo1985}.

\subsection{Quasi-linear (Fung) viscoelasticity}
With $\beta_i = g_i/(1-\sum_j g_j)$, the total resultants are
$\vect{f}_\mathrm{loc} = \vect{f}_\infty + \sum_i \vect{q}_i$, and each Prony branch is driven by increments
of the long-term resultants over a step of length $h$:
\begin{equation}
\vect{q}_i^{\,n+1} = a_i\,\vect{q}_i^{\,n} + \beta_i\,b_i\left(\vect{f}_\infty^{\,n+1}-\vect{f}_\infty^{\,n}\right),
\qquad a_i = e^{-h/\tau_i},\quad b_i = \frac{1-a_i}{h/\tau_i} = -\frac{\operatorname{expm1}(-h/\tau_i)}{h/\tau_i}.
\label{eq:prony}
\end{equation}
This is exact for a constant rate within the step. For $h\to0$, $b_i\to1$ and the response is instantaneous,
$\vect{f}_\mathrm{loc} = (E_0/E_\infty)\vect{f}_\infty$ to first order, which is the stiffness the stable time
step must use. The branch histories $(\vect{q}_i,\vect{f}_\infty^{\,n})$ are part of the state.

%=====================================================================================================
\section{Inertia}
\label{sec:mass}
Lumped translational mass per element end $\rho A L_0/2$. Rotary inertia per element end (scalar, isotropic)
\begin{equation}
j_e = \frac{\rho L_0}{2}\max\!\left(I_y + I_z,\ \frac{A L_0^2}{12}\right),
\end{equation}
the standard explicit rotary-inertia scaling (it only activates for $L_0 > \sqrt{6}\,r$, i.e.\ not on the
1\,mm results mesh). With isotropic nodal rotary inertia the rotational equation is $j\,\dot{\vect\omega} =
\vect{m}$ with no gyroscopic term.

%=====================================================================================================
\section{Test rig and boundary conditions}
\label{sec:rig}
The rig reproduces a drop-weight test on a coupon bonded between two plates:
\begin{itemize}
\item \textbf{Base plate (bonded, rigid, fixed).} All nodes on the bottom face ($z=0$) have all six degrees of
freedom fixed.
\item \textbf{Top plate (bonded, rigid).} All nodes on the top face are tied into one rigid body that
translates (rotations fixed). Its mass is the plate mass ($2$\,mm aluminium, $7.8$\,g) plus the lumped mass of
the tied nodes. The tie is imposed on accelerations: for tied set $\mathcal{T}$,
$\vect{a}_a \leftarrow \sum_{b\in\mathcal T} m_b\vect{a}_b / \sum_{b\in\mathcal T} m_b$ for all $a\in\mathcal T$,
so the plate moves under the resultant of all forces on its nodes.
\item \textbf{Impactor.} A rigid body with one translational degree of freedom ($z$), a separate node without
elements, interacting with the top plate only through contact (Section~\ref{sec:impactor-contact}).
\item \textbf{Loading: co-moving start.} At $t=0$ the impactor and the top plate move together at
$v_0 = 6$\,m/s towards the base, with an impactor--plate gap of $0.1$\,mm, and the impactor mass is set so
that the total kinetic energy equals the impact energy $E_\mathrm{imp}$ ($0.75$\,J):
\begin{equation}
M = \frac{2E_\mathrm{imp}}{v_0^2} - m_\mathrm{plate} .
\end{equation}
This takes the (plastic) impactor--plate collision as completed, i.e.\ the plate acts as the striker face.
The contact force then builds up continuously as the lattice decelerates the plate and releases continuously on
rebound. (Starting from rest with a gap gives a rigid--rigid collision, a force spike and rattling of the light
plate; this is physical but not what is measured.)
\item \textbf{Crush mode.} For constant-velocity crushing (the standard specific-energy test) the impactor mass
is set to $1$\,t, so the plate keeps $v_0$ independently of the lattice; the run stops at the target stroke.
\item Gravity is neglected ($\Delta v = g\,t \approx 0.1$\,m/s over 10\,ms).
\end{itemize}

%=====================================================================================================
\section{Contact and friction}
\label{sec:contact}
All normal contact laws use the smooth penetration
\begin{equation}
\delta(g) = \varepsilon\,\softplus(-g/\varepsilon) ,
\end{equation}
which tends to $\max(-g,0)$ for $|g|\gg\varepsilon$ and is $C^\infty$. Unless stated otherwise
$\varepsilon = 0.05\,r_{\min}$. Under light load contact therefore acts within $\approx2\varepsilon$ of
geometric touching (the regularisation width).

\subsection{Impactor--plate contact: critically damped Hertz law}
\label{sec:impactor-contact}
With gap $g = g_0 + u_{z,\mathrm{imp}} - u_{z,\mathrm{plate}}$ and gap rate $\dot g$, the contact force on the
impactor is
\begin{equation}
f = k\,\delta^{3/2} - c(\delta)\,\dot g,\qquad c(\delta) = 2\zeta\sqrt{k_t\,\mu_r},\quad k_t = \tfrac32 k\,\delta^{1/2},
\quad \mu_r = \frac{M\,m_\mathrm{plate}}{M+m_\mathrm{plate}},\quad \zeta = 1,
\end{equation}
followed by a smooth positive part that is exactly zero at $f=0$:
\begin{equation}
F = \max\!\left(0,\ \tfrac12\left(f + \sqrt{f^2 + s^2}\right) - \tfrac12 s\right),\qquad s = 0.01\,k\,\delta^{3/2}.
\end{equation}
$k$ is calibrated so that the law carries $F_\mathrm{ref} = 2000$\,N at $\delta_\mathrm{ref} = 0.2$\,mm
($k = F_\mathrm{ref}/\delta_\mathrm{ref}^{3/2}$, $\varepsilon = 0.1\,\delta_\mathrm{ref}$). The damping, tuned
to the impactor--plate contact mode itself, removes the ringing of the light plate between the heavy impactor
and the soft lattice. It vanishes as $\delta^{1/4}$ at onset (continuous force), and the positive part prevents
sticky tension at separation. Alternatives implemented for comparison: linear penalty, Hertz, Hunt--Crossley
\cite{hunt1975,lankarani1990} and an IPC log-barrier \cite{li2020ipc}.

\subsection{Beam--beam contact: line-contact density}
\label{sec:beam-beam}
Every element is a capsule: its centreline segment with radius $r_e$. Rather than one point force per pair of
segments, contact is a \emph{density integrated along the beams}.

\paragraph{Quadrature points.} Each element $a$ carries $n_q = \lceil L_{\max}/h_q\rceil$ points
$\vect{p}_{a,k} = \vect{x}_{a_1} + \xi_k(\vect{x}_{a_2}-\vect{x}_{a_1})$, $\xi_k = (k+\tfrac12)/n_q$, with
tributary length $w_a = L_{0,a}/n_q$ ($h_q = 0.1$\,mm).

\paragraph{Gap to another beam.} For a point $\vect p$ and an element $b$ with end points $\vect q_0$,
$\vect q_0 + \vect d_b$, the projection is
$t = \operatorname{clamp}\left((\vect p-\vect q_0)\cdot\vect d_b/(\vect d_b\cdot\vect d_b),\,0,\,1\right)$, and
the gap of point $(a,k)$ to beam $B$ (all candidate elements of $B$) is the distance to the polyline minus both
radii,
\begin{equation}
g_{a,k}^{B} = \min_{b\in B}\ \left\lVert \vect p_{a,k} - \vect q_0^b - t\,\vect d_b\right\rVert - r_a - r_B .
\end{equation}
Taking the minimum over the elements of $B$ (rather than summing over element pairs) avoids counting a contact
twice where two elements of $B$ share a node. $t$ is kept in the derivative: the force (first derivative) would
be exact with $t$ frozen (envelope theorem), but the contact \emph{stiffness}, needed for trajectory
sensitivities, contains $\partial t/\partial\vect x$.

\paragraph{Potential.}
\begin{equation}
\Phi_\mathrm{bb} = \frac12\sum_{a,k}\sum_{B\neq\mathrm{beam}(a)} w_a\,\frac{k_\ell}{2}\,\delta\!\left(g_{a,k}^{B}\right)^2,
\qquad k_\ell = \frac{E^*}{2},\quad E^* = \frac{E_0}{2(1-\nu^2)},
\end{equation}
and the nodal forces are $-\partial\Phi_\mathrm{bb}/\partial\vect x$ (automatic differentiation), so the
contact is conservative. The factor $\tfrac12$ accounts for each contact zone being integrated from both beams.

\paragraph{Line and point contact from one law.} Two parallel beams overlapping over a length $\ell$ carry
$F = k_\ell\,\delta\,\ell$ (line contact). For two crossed cylinders of radius $R$ the capsule gap along one
beam is $g(x) \approx -\delta + x^2/(4R)$, and integrating the density over the contact zone gives
\begin{equation}
F = \int k_\ell\left(\delta - \frac{x^2}{4R}\right)_{\!+}\dd x = \frac{4}{3}\,k_\ell\,\sqrt{4R}\,\delta^{3/2}
= \frac{4}{3}E^*\sqrt{R}\,\delta^{3/2}\quad\text{for } k_\ell = E^*/2,
\end{equation}
i.e.\ exactly the Hertz law for crossed cylinders~\cite{johnson1985}, without using it. (Numerically: 0.99,
0.98, 0.97 of Hertz at $\delta = 0.02$, $0.05$, $0.1$\,mm.)

\paragraph{Exclusions (geometric, mesh independent).}
\begin{enumerate}
\item A beam does not contact itself.
\item A (point, beam) combination already overlapping in the reference configuration ($g\le0$) never contacts.
\item Joint crotch: if the two beams share a lattice node, a (point, beam) combination whose reference gap is
below one radius is treated as joint material. Otherwise any joint rotation drives the crotch surfaces into
``contact'' and double counts the joint stiffness, which the beam element already provides.
\end{enumerate}
Contact between struts of the same joint further out (folding struts) remains active.

\paragraph{Candidate search.} Element pairs from different beams with segment--segment gap below a margin
($1$\,mm) are kept in a Verlet list. It is refreshed when any node has moved more than half the margin since
the last refresh, so no pair can close unseen. For each step only the pairs whose segment gap is below
$6\varepsilon$ are evaluated point by point; beyond this the potential is below $e^{-12}$ of the contact
scale. Both choices are deterministic functions of the state (needed for the adjoint).

\subsection{Beam--plate contact}
Every non-bonded lattice node $a$ contacts both plates with a Hertz potential of a sphere of the nodal radius
$r_a$ (largest radius of the attached elements) on a rigid flat:
\begin{equation}
\Phi_\mathrm{pl} = \sum_{a} \frac{2}{5}k_a\,\delta(g_a)^{5/2},\qquad k_a = \frac43\frac{E_0}{1-\nu^2}\sqrt{r_a},
\qquad g_a^\mathrm{base} = z_a - r_a,\quad g_a^\mathrm{top} = z_\mathrm{plate} - z_a - r_a .
\end{equation}
The reaction on the top plate enters the plate's equation of motion through $z_\mathrm{plate}$.

\subsection{Friction}
Coulomb friction with $\mu = 0.5$ acts at every contacting line-contact point and plate contact. With normal
force $F_n$ (from the same potential) and tangential slip velocity $\vect v_t$ between the two contacting
material points (the point on beam $a$ and its closest point on beam $b$, or node and plate),
\begin{equation}
\vect F_t = -\mu F_n\,\frac{\vect v_t}{\sqrt{\lVert\vect v_t\rVert^2 + v_\varepsilon^2}},\qquad
v_\varepsilon = \max\!\left(v_{\varepsilon 0},\ \frac{2\mu F_n h}{m_\mathrm{eff}}\right),\quad v_{\varepsilon0} = 1\,\mathrm{mm/s},
\end{equation}
with $m_\mathrm{eff} = m_p m_q/(m_p+m_q)$ the reduced mass of the two contact points (masses interpolated from
the element nodes). The equivalent viscous coefficient $\mu F_n/v_\varepsilon$ never exceeds
$m_\mathrm{eff}/(2h)$, so near stick the explicit step stays stable and the law acts as a stick constraint for
slip speeds below $\approx\mu F_n h/m_\mathrm{eff}$; under real slip it is Coulomb's $\mu F_n$. The forces are
applied with the interpolation weights at both contact points (equal and opposite).

%=====================================================================================================
\section{Explicit time integration}
\label{sec:time}
\subsection{Velocity Verlet with exponential-map rotations}
With nodal accelerations $\vect a = (\vect f_\mathrm{ext}+\vect f_\mathrm{c}-\vect f_u)/m$ and angular
accelerations $\vect\alpha = -\vect f_\theta/j$ (fixed degrees of freedom masked), one step of length $h_n$ is
\begin{align}
\vect v^{n+1/2} &= \vect v^n + \tfrac{h_n}{2}\vect a^n, & \vect\omega^{n+1/2} &= \vect\omega^n + \tfrac{h_n}{2}\vect\alpha^n,\\
\vect u^{n+1} &= \vect u^n + h_n\vect v^{n+1/2}, & \mat R^{n+1} &= \exp\!\left(\skw{h_n\vect\omega^{n+1/2}}\right)\mat R^n,\\
\vect a^{n+1},\ \vect\alpha^{n+1} &= \text{forces at } (\vect u^{n+1}, \mat R^{n+1}, \vect v^{n+1/2}), &&\\
\vect v^{n+1} &= \vect v^{n+1/2} + \tfrac{h_n}{2}\vect a^{n+1}, & \vect\omega^{n+1} &= \vect\omega^{n+1/2} + \tfrac{h_n}{2}\vect\alpha^{n+1},
\end{align}
followed by the Prony update (\ref{eq:prony}). The exponential map is evaluated with the Rodrigues formula in a
form that is smooth at zero rotation,
$\exp(\skw{\vect w}) = \mat I + \operatorname{sinc}(r)\skw{\vect w} + \tfrac12\operatorname{sinc}(r/2)^2\skw{\vect w}^2$,
$r = \lVert\vect w\rVert$. One force evaluation per step; velocity-dependent forces (damping, friction) see
the half-step velocity.

\subsection{Stable time step}
The step is adapted every step as
\begin{equation}
h_n = \min\!\left(\frac{2C}{\sqrt{\omega_\mathrm{beam}^2 + \omega_\mathrm{c}^2}},\ h_\mathrm{cap}\right),\qquad C = 0.7,
\end{equation}
where:
\begin{itemize}
\item $\omega_\mathrm{beam}^2 = \lambda_{\max}(\mat M^{-1/2}\mat K\mat M^{-1/2})$ of the current tangent is
estimated by a warm-started power iteration (Jacobian--vector products of the internal force, free--free,
i.e.\ an upper bound for any support), refreshed every 20 steps. With $h\ll\tau_1$ the Prony law makes the
tangent the instantaneous ($E_0$) one.
\item $\omega_\mathrm{c}$ collects contact stiffnesses. For the impactor law,
$\omega = \sqrt{k_t/\mu_r}/(\sqrt{1+\zeta^2}-\zeta)$ (damped central-difference limit). For line and plate
contact, every quadrature point's exact stiffness
$\tfrac12 w_a k_\ell(\sigma^2 + \delta\sigma(1-\sigma)/\varepsilon)$ is summed per element and per node (a
row-sum bound, which covers nodes in several contacts at once), giving $\sqrt{2K_\mathrm{node}/m_\mathrm{node}}$.
\item $h_\mathrm{cap}$ is a kinematic cap: near contact, the relative approach per step is limited to
$0.1\ell$ ($\ell = \delta_\mathrm{ref}$ for the impactor law, $2\varepsilon$ for lattice contact), so the force
cannot jump from zero to a large value within one step. The stiffness-based estimate cannot see this, because
the stiffness is zero just before contact.
\end{itemize}
If $h_n$ falls below $10^{-3}$ of its initial value the run stops (instability guard). Typical values: $h\approx
4\times10^{-6}$\,s on the 1\,mm mesh with $E_0$; an 8\,ms drop takes $\approx 7400$ steps.

%=====================================================================================================
\section{Energy accounting}
\label{sec:energy}
The energy balance is monitored as
\begin{multline}
E_\mathrm{tot} = \underbrace{\tfrac12\sum m\lVert\vect v\rVert^2 + \tfrac12\sum j\lVert\vect\omega\rVert^2}_{\text{kinetic}}\\
+ \underbrace{\sum_n \tfrac12(\vect f_\mathrm{loc}^n + \vect f_\mathrm{loc}^{n+1})\cdot(\vect d^{n+1}-\vect d^n)}_{\text{internal work (stored + dissipated)}}\\
+ \Phi_\mathrm{imp} + W_\mathrm{damp,imp} + \Phi_\mathrm{bb} + \Phi_\mathrm{pl} + W_\mathrm{fric},
\end{multline}
where $W_\mathrm{fric} = -\sum_n \tfrac12(\vect F_t^n+\vect F_t^{n+1})\cdot(\vect u^{n+1}-\vect u^n)$ is the
work of the friction forces actually applied. The material dissipation is the internal work minus the stored
long-term energy (fibre Yeoh plus quadratic parts).

%=====================================================================================================
\section{Response measures, objectives and exact sensitivities}
\label{sec:sens}
\subsection{Response curve and objectives}
The response is the impactor force $F$ versus top-plate stroke $\delta$. Per step,
$\Delta\delta = \delta^{n+1}-\delta^n$, $\bar F = \tfrac12(F^n+F^{n+1})$, $\bar\delta = \tfrac12(\delta^n+\delta^{n+1})$, and
\begin{equation}
S_0 = \sum_n \left(\bar F - F^*(\bar\delta)\right)^2\relu(\Delta\delta),\qquad
S_1 = \sum_n \bar F\,\relu(\Delta\delta),\qquad S_2 = \sum_n \bar F\,\Delta\delta .
\end{equation}
\begin{itemize}
\item \textbf{Target force--displacement curve (inverse design):}
$J_\mathrm{MSE} = S_0 \big/ \int F^{*2}\dd\delta$, with $F^*$ piecewise linear (e.g.\ an ``ideal absorber''
plateau~\cite{zeng2023}).
\item \textbf{Specific energy absorption:} $\mathrm{SEA} = S_k/m_\mathrm{lat}$ with
$m_\mathrm{lat} = \rho\sum A L_0$; $k=1$ (loading area) for constant-velocity crushing, $k=2$ (dissipated loop
area) in a fixed-energy drop. There the loading area is close to the impact energy whatever the design, so the
loading-area SEA would reduce to $1/m$. The objective $J = -\mathrm{SEA}$ is minimised.
\end{itemize}

\subsection{Discrete adjoint with two-level checkpointing}
Write one time step as the pure map $\vect z^{n+1} = \Phi_n(\vect z^n;\vect\theta)$ with state
$\vect z = (\vect u,\vect v,\mat R,\vect\omega,\vect a,\vect\alpha,\vect q,\vect f_\infty,F_\mathrm{b},F)$, and
the objective as $J = g(\vect S,\vect\theta)$ with additive accumulators
$\vect S = \sum_n\vect\ell_n(\vect z^n,\vect z^{n+1})$. The time-step schedule $\{h_n\}$ of the nominal run is
replayed, and design-dependent rig constants (impactor mass and law, $\varepsilon$, exclusion sets, plate
contact node sets) are frozen at the nominal design. The contact Verlet list is stored with each checkpoint and
restored on recompute. The adjoint recursion is
\begin{equation}
\vect\lambda^n = \left(\frac{\partial\Phi_n}{\partial\vect z^n}\right)^{\!\top}\vect\lambda^{n+1}
+ \frac{\partial (\vect w\cdot\vect\ell_n)}{\partial\vect z^n},\qquad
\frac{\dd J}{\dd\vect\theta} = \frac{\partial g}{\partial\vect\theta}
+ \sum_n\left[\left(\frac{\partial\Phi_n}{\partial\vect\theta}\right)^{\!\top}\vect\lambda^{n+1}
+ \frac{\partial(\vect w\cdot\vect\ell_n)}{\partial\vect\theta}\right],\qquad \vect w = \frac{\partial g}{\partial\vect S},
\end{equation}
with each product a single-step vector--Jacobian product.

\begin{algorithm}[h]
\caption{Checkpointed discrete adjoint (\texttt{adjoint.checkpointed\_adjoint})}
\begin{algorithmic}[1]
\State forward (no graph): $\vect z^0$; for $n=0..N-1$: store $\vect z^n$ and contact state if $n \bmod K = 0$; $\vect z^{n+1}=\Phi_n(\vect z^n)$, $\vect S \mathrel{+}= \vect\ell_n$
\State $\vect w \gets \partial g/\partial\vect S$, $\vect g_\theta \gets \partial g/\partial\vect\theta$, $\vect\lambda\gets 0$
\For{segments $[s, s+K)$ from last to first}
  \State restore checkpoint $s$; recompute and store $\vect z^{s..s+K-1}$ (no graph)
  \For{$n = s+K-1$ down to $s$}
    \State build the graph of one step $\vect z^{n+1},\vect\ell_n = \Phi_n(\vect z^n)$ from the stored $\vect z^n$
    \State $(\vect\lambda, \Delta\vect g_\theta) \gets \mathrm{VJP}\big((\vect z^{n+1}, \vect w\cdot\vect\ell_n);\ (\vect\lambda, 1)\big)$; $\vect g_\theta \mathrel{+}= \Delta\vect g_\theta$
  \EndFor
\EndFor
\State chain $\vect\lambda^0$ through $\vect z^0(\vect\theta)$
\end{algorithmic}
\end{algorithm}

Memory is $\mathcal O((N/K + K)\,|\vect z|)$ plus one step graph, $K=\sqrt N$, independent of the number of
design variables. The gradient equals reverse-mode differentiation of the unrolled loop
(discretise-then-differentiate~\cite{jensen2014}); checkpointing follows~\cite{griewank2000}. Measured on the
3\,mm mesh at 1\,ms: unrolled reverse AD 327\,s and 16.7\,GB (CPU), adjoint 29.8\,s and 0.69\,GB; the full
8\,ms drop on the 1\,mm mesh takes 403\,s and 0.43\,GB on an 8\,GB GPU.

\subsection{Optimisers}
Design variables are normalised to $[0,1]$ and denormalised inside the objective. For many design
variables with constraints the method of moving asymptotes (MMA,~\cite{svanberg1987}) is used with asymptotes
carried between iterations, or its globally convergent variant GCMMA~\cite{svanberg2002}, whose inner loop
enlarges the approximation curvature using value-only (forward) evaluations until the approximation is
conservative. L-BFGS-B is used for small unconstrained problems.

%=====================================================================================================
\section{Verification}
\label{sec:verification}
\begin{longtable}{p{0.42\textwidth}p{0.5\textwidth}}
\toprule
check & result\\
\midrule
small-strain tangent vs.\ classical Timoshenko stiffness & equal to $10^{-9}$ (skewed and bent frames, $I_y\neq I_z$)\\
rigid rotations up to $172^\circ$ & internal force $<10^{-10}$\\
geometric stiffness: column under axial load $P$ & $\omega_1^2(P)/\omega_1^2(0) = 1-P/P_\mathrm{cr}$ within 2\,\% (Engesser $P_\mathrm{cr}$)\\
cantilever first frequency & within 0.12\,\% of Euler--Bernoulli\\
tip moment $2\pi EI/L$ & beam closes to a full circle (tip error $1.4\times10^{-5}$\,mm)\\
fibre Yeoh: bending moment, torque & within 0.1\,\% / 1\,\% of direct section integration\\
Prony: step relaxation; loss per cycle at 1\,kHz & exact $E_\infty + (E_0-E_\infty)e^{-t/\tau}$; $\pi E''\varepsilon_0^2V$ within 2\,\%\\
line contact, crossed cylinders & Hertz within 1--3\,\% (from integration alone)\\
beam radius in contact & no force for $d > r_A + r_B$; onset at $r_A+r_B$ (radii 0.8/0.5)\\
two free beams colliding & energy within 1\,\%, momentum to $10^{-9}$, no interpenetration\\
friction & $\lVert\vect F_t\rVert = \mu F_n$ under slip; power $=\mu F_n\lVert\vect v\rVert$; equal and opposite\\
drop energy balance & $\le 1.1\times10^{-4}$ relative (all physics)\\
contact force continuity & per-step force change $\propto h$ (no jump)\\
mesh convergence (0.75 vs 1.0\,mm) & stroke 2\,\%, dissipation 0.2\,\%, densification peak $\pm3.5$\,\%\\
adjoint vs.\ unrolled AD; vs.\ finite differences & $10^{-16}$; $10^{-9}$--$10^{-10}$ (non-uniform designs, all physics)\\
\bottomrule
\end{longtable}

%=====================================================================================================
\section{Limitations}
\label{sec:limits}
\begin{itemize}
\item One Prony term ($\tau_1 = 52\,\mu$s): most of the dissipation occurs within $\approx 1$\,ms, and the slow
viscoelasticity near $T_g$ is absent. The relaxation spectrum is the most influential uncertain input.
\item Transverse shear is linear; only the axial, bending and torsional fibre responses are hyperelastic.
\item Joints are centreline intersections: overlapping material is double counted in the mass, and joint
stiffness comes from the beam ends only.
\item Contact acts within $\approx2\varepsilon$ of touching under light load (regularisation width
$\varepsilon = 0.05\,r_{\min}$).
\item The co-moving start assumes a completed plastic impactor--plate collision.
\item Friction mostly sticks at these loads and dissipates $<1$\,\% of the impact energy.
\end{itemize}

%=====================================================================================================
\begin{thebibliography}{99}
\bibitem{battini2002} J.-M. Battini, C. Pacoste. Co-rotational beam elements with warping effects in instability problems. \emph{Comput. Methods Appl. Mech. Eng.} 191 (2002) 1755--1789.
\bibitem{simo1985} J.C. Simo. A finite strain beam formulation. The three-dimensional dynamic problem. Part I. \emph{Comput. Methods Appl. Mech. Eng.} 49 (1985) 55--70.
\bibitem{cowper1966} G.R. Cowper. The shear coefficient in Timoshenko's beam theory. \emph{J. Appl. Mech.} 33 (1966) 335--340.
\bibitem{yeoh1993} O.H. Yeoh. Some forms of the strain energy function for rubber. \emph{Rubber Chem. Technol.} 66 (1993) 754--771.
\bibitem{fung1993} Y.C. Fung. \emph{Biomechanics: Mechanical Properties of Living Tissues}. Springer, 1993 (quasi-linear viscoelasticity).
\bibitem{johnson1985} K.L. Johnson. \emph{Contact Mechanics}. Cambridge University Press, 1985.
\bibitem{hunt1975} K.H. Hunt, F.R.E. Crossley. Coefficient of restitution interpreted as damping in vibroimpact. \emph{J. Appl. Mech.} 42 (1975) 440--445.
\bibitem{lankarani1990} H.M. Lankarani, P.E. Nikravesh. A contact force model with hysteresis damping for impact analysis of multibody systems. \emph{J. Mech. Des.} 112 (1990) 369--376.
\bibitem{li2020ipc} M. Li et al. Incremental potential contact. \emph{ACM Trans. Graph.} 39 (2020) 49.
\bibitem{svanberg1987} K. Svanberg. The method of moving asymptotes---a new method for structural optimization. \emph{Int. J. Numer. Methods Eng.} 24 (1987) 359--373.
\bibitem{svanberg2002} K. Svanberg. A class of globally convergent optimization methods based on conservative convex separable approximations. \emph{SIAM J. Optim.} 12 (2002) 555--573.
\bibitem{griewank2000} A. Griewank, A. Walther. Algorithm 799: Revolve: an implementation of checkpointing for the reverse or adjoint mode of computational differentiation. \emph{ACM Trans. Math. Softw.} 26 (2000) 19--45.
\bibitem{jensen2014} J.S. Jensen, P.B. Nakshatrala, D.A. Tortorelli. On the consistency of adjoint sensitivity analysis for structural optimization of linear dynamic problems. \emph{Struct. Multidiscip. Optim.} 49 (2014) 831--837.
\bibitem{zeng2023} Q. Zeng et al. Inverse design of energy-absorbing metamaterials by topology optimization. \emph{Adv. Sci.} 10 (2023) 2204977.
\bibitem{lebarbenchon2024} L. Le Barbenchon, M. Li{\ss}ner. On the dynamic performance of additively manufactured visco-elastic meta-materials. \emph{Mater. Lett.} 359 (2024) 135823.
\bibitem{xu2025} K. Xu, J. Ji, H. Feng, J. Lin, Y. Li. Multifunctional performance of 3D-printed helmet liners under impact and blast loadings. \emph{Thin-Walled Struct.} 218 (2025) 114036.
\bibitem{carbon} Carbon Inc. EPU 45 technical data sheet.
\end{thebibliography}

\end{document}
